\startslide
\commonframe
\mbox{}

\begin{tikzpicture}[remember picture,overlay]
  \slidetitle{Ключевые результаты}

  \node[
    anchor=north west,
    fill=softbg,
    rounded corners=10pt,
    text width=4.20cm,
    inner xsep=0.28cm,
    inner ysep=0.25cm,
    align=left
  ] at ([xshift=0.65cm,yshift=-1.95cm]current page.north west) {%
    {\fontsize{9.0}{10.2}\selectfont
    \textbf{PostgreSQL}\\[0.10cm]
    Датасет: \texttt{1,000,000} операций.\\
    Силен в точечных и умеренно\\
    селективных запросах.}%
  };

  \node[
    anchor=north west,
    fill=softrose,
    rounded corners=10pt,
    text width=4.20cm,
    inner xsep=0.28cm,
    inner ysep=0.25cm,
    align=left
  ] at ([xshift=5.58cm,yshift=-1.95cm]current page.north west) {%
    {\fontsize{9.0}{10.2}\selectfont
    \textbf{MongoDB}\\[0.10cm]
    \texttt{10,000} документов и около \texttt{5.08} млн встроенных операций.\\
    Лучше всего работает с полной историей объекта.}%
  };

  \node[
    anchor=north west,
    fill=softcyan,
    rounded corners=10pt,
    text width=4.20cm,
    inner xsep=0.28cm,
    inner ysep=0.25cm,
    align=left
  ] at ([xshift=10.51cm,yshift=-1.95cm]current page.north west) {%
    {\fontsize{9.0}{10.2}\selectfont
    \textbf{YTsaurus}\\[0.10cm]
    Лучшие времена у запросов по префиксу ключа:\\
    \texttt{136 ms} для поиска по\\
    \texttt{objectId + replicaId + txId}.}%
  };

  \node[anchor=north west, text=textblue] at ([xshift=1.00cm,yshift=-4.85cm]current page.north west) {%
    {\fontsize{11.0}{12.6}\selectfont\bfseries Что показало сравнение}%
  };

  \node[anchor=north west, align=left, text width=13.7cm] at ([xshift=1.05cm,yshift=-5.72cm]current page.north west) {%
    {\fontsize{9.0}{10.2}\selectfont
    • \textbf{PostgreSQL} естественен для глобального immutable operation log.\\
    • \textbf{MongoDB} удобен для объектно-центричной истории и вложенных структур.\\
    • \textbf{YTsaurus} выигрывает там, где запрос согласован с составным ключом.}%
  };
\end{tikzpicture}

\footerband{7}
